\section{RQ4: Diagnosis Actionability}
The preregistered human study contains 40 stratified failure cases and compares
a semantic score/claim output against the full evidence-chain report. Option
order is counterbalanced and condition names are hidden. At least three
non-author reviewers must rate diagnosis correctness, repair actionability,
specificity, evidence sufficiency, preference, confidence, and time. No human
responses are currently available; the primary human RQ4 outcome is therefore
N/A.

We separately execute a controlled LLM-simulated protocol stress test using
three role prompts: strict evidence, developer actionability, and skeptical ARR
review. Each role evaluates all 40 cases under raw trace (A), score-only (B), and
evidence-chain (C) settings, yielding 360 RQ4 judgments. The model is pinned to
\texttt{gpt-4.1-nano-2025-04-14}; strict schema validation produces zero parse
failures. Private labels are loaded only after responses are frozen.

Table~\ref{tab:table4-rq4-actionability} shows the simulated results. Setting C
raises normalized root-cause accuracy from 0.067/0.117 (A/B) to 0.908 and lowers
false greens from 0.625/0.550 to 0.400. However, mean actionability falls from
3.767 (A) and 3.608 (B) to 3.333 (C). Thus the evidence-chain condition does not
improve every registered proxy, and the preregistered positive simulated-pilot
claim is not supported. The fix-correctness value is an operational proxy, not
a human rating. Figure~\ref{fig:rq4} summarizes the protocol.

We include this controlled LLM-simulated review and actionability pilot as a
stress test of the review protocol. These results are reported separately and
do not replace independent human validation.

% Generated by paper/generate_tables.py; do not edit manually.
\begin{table*}[t]
\centering
\scriptsize
\resizebox{\textwidth}{!}{%
\begin{tabular}{llllllll}
\toprule
 Study & Review status & Cases & Reviewers / agents & Root cause & Fix proxy & Actionability & False green \\
\midrule
 Preregistered paired actionability study & pending three independent reviewers & 40 & 0/3 humans & N/A & N/A & N/A & N/A \\
 Simulated A: raw trace & LLM-simulated pilot; not human validation & 40 & 3 & 0.067 & 0.008 & 3.767 & 0.625 \\
 Simulated B: score only & LLM-simulated pilot; not human validation & 40 & 3 & 0.117 & 0.017 & 3.608 & 0.550 \\
 Simulated C: evidence chain & LLM-simulated pilot; not human validation & 40 & 3 & 0.908 & 0.342 & 3.333 & 0.400 \\
\bottomrule
\end{tabular}%
}
\caption{Table4 Rq4 Actionability. Review status and unavailable measurements are shown explicitly.}
\label{tab:table4-rq4-actionability}
\end{table*}

\begin{figure}[t]
\centering
\begin{tikzpicture}[
  node distance=5mm,
  item/.style={draw, rounded corners=2pt, align=left, text width=68mm, font=\footnotesize, inner sep=4pt},
  arrow/.style={-{Latex[length=1.8mm]}, thick}
]
\node[item] (case) {Fixed failure case: query, answer, and retrieved contexts};
\node[item, below=of case] (options) {Blinded option order: semantic-core output versus full evidence-chain output};
\node[item, below=of options] (ratings) {Reviewer judgments: correctness, actionability, specificity, evidence sufficiency, time};
\node[item, below=of ratings] (analysis) {Preregistered analysis: decisive preference, paired deltas, intervals, and disagreement};
\draw[arrow] (case) -- (options);
\draw[arrow] (options) -- (ratings);
\draw[arrow] (ratings) -- (analysis);
\end{tikzpicture}
\caption{RQ4 study settings. The protocol is frozen; responses have not yet
been collected.}
\label{fig:rq4}
\end{figure}

